\documentclass[a4paper,12pt]{article}

\usepackage[a4paper,top=20mm,bottom=24mm,left=18mm,right=18mm]{geometry}
\usepackage{iftex}
\ifXeTeX
  \usepackage{fontspec}
  \usepackage{polyglossia}
  \setdefaultlanguage{russian}
  \setotherlanguage{english}
  \setmainfont{Times New Roman}
  \setmonofont{Menlo}
  \newfontfamily\cyrillicfont{Times New Roman}
  \newfontfamily\cyrillicfonttt{Menlo}
\else
  \usepackage[T1,T2A]{fontenc}
  \usepackage[utf8]{inputenc}
  \usepackage[english,russian]{babel}
\fi
\usepackage{xcolor}
\usepackage{enumitem}
\usepackage{hyperref}
\usepackage{listings}
\usepackage{titlesec}

\hypersetup{colorlinks=true,urlcolor=blue,linkcolor=black}
\setlength{\parindent}{0pt}
\setlength{\parskip}{6pt}
\linespread{1.12}
\sloppy

\definecolor{codebg}{RGB}{247,248,250}
\definecolor{codeframe}{RGB}{214,219,226}
\definecolor{codekw}{RGB}{20,78,140}

\lstdefinestyle{cmd}{
  basicstyle=\small\ttfamily,
  backgroundcolor=\color{codebg},
  frame=single,
  rulecolor=\color{codeframe},
  breaklines=true,
  columns=fullflexible,
  keepspaces=true,
  showstringspaces=false,
  keywordstyle=\color{codekw},
}
\lstset{style=cmd}

\titleformat{\section}{\large\bfseries}{\thesection.}{0.6em}{}
\titleformat{\subsection}{\normalsize\bfseries}{\thesubsection.}{0.6em}{}

\title{\textbf{Команды проекта для расчётов и визуализаций}}
\author{Проект: точки Лагранжа в ограниченной задаче трёх тел}
\date{}

\begin{document}
\maketitle

\section*{Подготовка}

Все команды ниже выполняются из корня склонированного репозитория. То есть
сначала нужно перейти в папку проекта, а затем активировать виртуальное
окружение:

\begin{lstlisting}
source .venv/bin/activate
\end{lstlisting}

Команды можно запускать двумя способами. Через модуль Python:

\begin{lstlisting}
python -m mechanics.cli.main ...
\end{lstlisting}

Или через установленную консольную команду:

\begin{lstlisting}
mechanics ...
\end{lstlisting}

\section{Главная команда для всех готовых материалов}

\begin{lstlisting}
python -m mechanics.cli.main showcase --system sun-earth
\end{lstlisting}

Эквивалентная короткая форма:

\begin{lstlisting}
mechanics showcase --system sun-earth
\end{lstlisting}

Эта команда пересобирает основной набор результатов проекта. Она:

\begin{itemize}[leftmargin=1.2cm]
  \item задаёт выбранную систему двух массивных тел;
  \item находит все точки Лагранжа $L_1,\ldots,L_5$;
  \item строит карту точек Лагранжа;
  \item считает базовую неустойчивую траекторию около $L_1$;
  \item считает базовую устойчивую траекторию около $L_4$;
  \item делает сравнение движения около $L_1$ и $L_4$;
  \item делает сравнение при одинаковом начальном возмущении;
  \item строит серию расчётов по амплитудам $dx$;
  \item строит серию расчётов по шагу интегрирования $dt$;
  \item собирает общий интерактивный HTML-файл.
\end{itemize}

Основные создаваемые файлы:

\begin{lstlisting}
results/sun_earth_lagrange_points_coordinates.json
results/runs/sun_earth_l1_unstable.json
results/runs/sun_earth_l1_unstable_trajectory.csv
results/runs/sun_earth_l4_stable.json
results/runs/sun_earth_l4_stable_trajectory.csv
results/figures/sun_earth_lagrange_points_map.png
results/figures/sun_earth_lagrange_points_map.svg
results/figures/sun_earth_l1_unstable_2d.png
results/figures/sun_earth_l1_unstable_2d.svg
results/figures/sun_earth_l1_unstable_3d.png
results/figures/sun_earth_l1_unstable_3d.svg
results/figures/sun_earth_l4_stable_2d.png
results/figures/sun_earth_l4_stable_2d.svg
results/figures/sun_earth_l4_stable_3d.png
results/figures/sun_earth_l4_stable_3d.svg
results/html/sun_earth_showcase.html
\end{lstlisting}

\section{Один расчёт около $L_1$}

\begin{lstlisting}
python -m mechanics.cli.main simulate --system sun-earth --point L1 --dx 1e-4 --t-max 20 --dt 0.002 --output-stem sun_earth_l1_unstable
\end{lstlisting}

Смысл параметров:

\begin{itemize}[leftmargin=1.2cm]
  \item \texttt{simulate} -- запустить один расчёт;
  \item \texttt{--system sun-earth} -- выбрать физическую систему Солнце--Земля;
  \item \texttt{--point L1} -- считать движение около точки Лагранжа $L_1$;
  \item \texttt{--dx 1e-4} -- начальное смещение по оси $x$;
  \item \texttt{--t-max 20} -- считать до времени 20 в нормированных единицах;
  \item \texttt{--dt 0.002} -- шаг вывода данных;
  \item \texttt{--output-stem sun\_earth\_l1\_unstable} -- базовое имя файлов результата.
\end{itemize}

Создаёт:

\begin{lstlisting}
results/runs/sun_earth_l1_unstable.json
results/runs/sun_earth_l1_unstable_trajectory.csv
results/figures/sun_earth_l1_unstable_2d.png
results/figures/sun_earth_l1_unstable_3d.png
results/html/sun_earth_l1_unstable.html
\end{lstlisting}

\section{Один расчёт около $L_4$}

\begin{lstlisting}
python -m mechanics.cli.main simulate --system sun-earth --point L4 --dx 1e-4 --vx 0.001 --vy 0.001 --t-max 12 --dt 0.002 --output-stem sun_earth_l4_stable
\end{lstlisting}

Отличия от расчёта около $L_1$:

\begin{itemize}[leftmargin=1.2cm]
  \item \texttt{--point L4} -- движение около треугольной точки $L_4$;
  \item \texttt{--vx 0.001} -- начальная скорость по $x$;
  \item \texttt{--vy 0.001} -- начальная скорость по $y$.
\end{itemize}

Создаёт:

\begin{lstlisting}
results/runs/sun_earth_l4_stable.json
results/runs/sun_earth_l4_stable_trajectory.csv
results/figures/sun_earth_l4_stable_2d.png
results/figures/sun_earth_l4_stable_3d.png
results/html/sun_earth_l4_stable.html
\end{lstlisting}

\section{Расчёт около любой точки Лагранжа}

\begin{lstlisting}
python -m mechanics.cli.main simulate --system sun-earth --point L2 --dx 1e-4 --t-max 10 --dt 0.002 --output-stem sun_earth_l2_example
\end{lstlisting}

В параметре \texttt{--point} можно указать:

\begin{lstlisting}
--point L1
--point L2
--point L3
--point L4
--point L5
\end{lstlisting}

В параметре \texttt{--system} можно выбрать систему:

\begin{lstlisting}
--system sun-earth    система Солнце-Земля
--system earth-moon   система Земля-Луна
--system sun-jupiter  система Солнце-Юпитер
\end{lstlisting}

\section{Расчёт с другим начальным смещением}

\begin{lstlisting}
python -m mechanics.cli.main simulate --system sun-earth --point L1 --dx 1e-5 --t-max 12 --dt 0.002 --output-stem l1_dx_1e_minus_5
\end{lstlisting}

Параметры начального смещения:

\begin{itemize}[leftmargin=1.2cm]
  \item \texttt{--dx} -- смещение по $x$;
  \item \texttt{--dy} -- смещение по $y$;
  \item \texttt{--dz} -- смещение по $z$.
\end{itemize}

Пример трёхмерного возмущения:

\begin{lstlisting}
python -m mechanics.cli.main simulate --system sun-earth --point L4 --dx 1e-4 --dy 1e-4 --dz 1e-4 --t-max 12 --dt 0.002 --output-stem l4_3d_offset
\end{lstlisting}

\section{Расчёт с начальной скоростью}

\begin{lstlisting}
python -m mechanics.cli.main simulate --system sun-earth --point L4 --dx 1e-4 --vx 0.001 --vy 0.001 --vz 0.0 --t-max 12 --dt 0.002 --output-stem l4_with_velocity
\end{lstlisting}

Параметры начальной скорости:

\begin{itemize}[leftmargin=1.2cm]
  \item \texttt{--vx} -- начальная скорость по $x$;
  \item \texttt{--vy} -- начальная скорость по $y$;
  \item \texttt{--vz} -- начальная скорость по $z$.
\end{itemize}

\section{Проверка влияния шага $dt$}

Более грубый шаг:

\begin{lstlisting}
python -m mechanics.cli.main simulate --system sun-earth --point L1 --dx 1e-4 --t-max 12 --dt 0.004 --output-stem l1_dt_0004
\end{lstlisting}

Более точный шаг:

\begin{lstlisting}
python -m mechanics.cli.main simulate --system sun-earth --point L1 --dx 1e-4 --t-max 12 --dt 0.001 --output-stem l1_dt_0001
\end{lstlisting}

Смысл проверки:

\begin{itemize}[leftmargin=1.2cm]
  \item большой \texttt{dt} даёт меньше точек, считается быстрее, но грубее;
  \item маленький \texttt{dt} даёт больше точек, считается медленнее, но точнее;
  \item в проекте это нужно, чтобы показать: уход от $L_1$ не является ошибкой численного шага.
\end{itemize}

\section{Параметр \texttt{sample-stride} для HTML}

\begin{lstlisting}
python -m mechanics.cli.main simulate --system sun-earth --point L1 --dx 1e-4 --t-max 20 --dt 0.002 --sample-stride 10 --output-stem l1_sampled
\end{lstlisting}

Параметр \texttt{--sample-stride 10} означает, что в интерактивные графики
попадёт каждая десятая точка траектории. Меньшее значение делает HTML подробнее,
но тяжелее. Большее значение делает HTML легче, но траектория становится менее
детальной.

\section{Как создать и открыть анимацию}

Отдельной команды только для анимации в проекте нет. Анимация создаётся как
часть интерактивного HTML-файла.

Пример для траектории около $L_1$:

\begin{lstlisting}
mechanics simulate --system sun-earth --point L1 --dx 1e-4 --t-max 20 --dt 0.002 --output-stem l1_animation
\end{lstlisting}

После выполнения команды появится файл:

\begin{lstlisting}
results/html/l1_animation.html
\end{lstlisting}

Открыть его в браузере на macOS можно так:

\begin{lstlisting}
open results/html/l1_animation.html
\end{lstlisting}

Пример для траектории около $L_4$:

\begin{lstlisting}
mechanics simulate --system sun-earth --point L4 --dx 1e-4 --vx 0.001 --vy 0.001 --t-max 12 --dt 0.002 --output-stem l4_animation
\end{lstlisting}

Открытие результата:

\begin{lstlisting}
open results/html/l4_animation.html
\end{lstlisting}

Для общего демонстрационного HTML:

\begin{lstlisting}
mechanics showcase --system sun-earth
open results/html/sun_earth_showcase.html
\end{lstlisting}

Если открывать через Visual Studio Code, можно найти файл в дереве проекта,
нажать по нему правой кнопкой и выбрать открытие в браузере. Второй вариант:
выбрать \texttt{Reveal in Finder}, а затем открыть файл двойным кликом.

\section{Сборка основного PDF-отчёта}

\begin{lstlisting}
tectonic docs/tex/main.tex --outdir docs/pdf
\end{lstlisting}

Команда создаёт или обновляет файл:

\begin{lstlisting}
docs/pdf/main.pdf
\end{lstlisting}

\section{Главные параметры команды \texttt{simulate}}

\begin{lstlisting}
--system         система: sun-earth, earth-moon, sun-jupiter
--point          точка Лагранжа: L1, L2, L3, L4, L5
--dx             начальное смещение по x
--dy             начальное смещение по y
--dz             начальное смещение по z
--vx             начальная скорость по x
--vy             начальная скорость по y
--vz             начальная скорость по z
--t-max          конечное время расчёта в нормированных единицах
--dt             шаг вывода траектории
--sample-stride  прореживание точек для интерактивного HTML
--output-stem    имя выходных файлов
\end{lstlisting}

\section*{Главная практическая команда}

Самая полезная команда для проекта:

\begin{lstlisting}
python -m mechanics.cli.main showcase --system sun-earth
\end{lstlisting}

Она пересобирает основные картинки, таблицы расчётов и общий интерактивный
HTML-файл.

\section*{Универсальные команды всех типов}

Подготовка окружения:

\begin{lstlisting}
source .venv/bin/activate
\end{lstlisting}

Пересобрать все основные результаты проекта:

\begin{lstlisting}
mechanics showcase --system sun-earth
\end{lstlisting}

Запустить один расчёт:

\begin{lstlisting}
mechanics simulate --system sun-earth --point L1 --dx 1e-4 --t-max 20 --dt 0.002 --output-stem my_run
\end{lstlisting}

Открыть HTML с анимацией одного расчёта:

\begin{lstlisting}
open results/html/my_run.html
\end{lstlisting}

Создать и открыть анимацию для $L_1$:

\begin{lstlisting}
mechanics simulate --system sun-earth --point L1 --dx 1e-4 --t-max 20 --dt 0.002 --output-stem l1_animation
open results/html/l1_animation.html
\end{lstlisting}

Создать и открыть анимацию для $L_4$:

\begin{lstlisting}
mechanics simulate --system sun-earth --point L4 --dx 1e-4 --vx 0.001 --vy 0.001 --t-max 12 --dt 0.002 --output-stem l4_animation
open results/html/l4_animation.html
\end{lstlisting}

Открыть общий интерактивный HTML:

\begin{lstlisting}
open results/html/sun_earth_showcase.html
\end{lstlisting}

Пересобрать основной PDF-отчёт:

\begin{lstlisting}
tectonic docs/tex/main.tex --outdir docs/pdf
\end{lstlisting}

Пересобрать этот PDF со списком команд:

\begin{lstlisting}
tectonic docs/tex/commands.tex --outdir docs/pdf
\end{lstlisting}

\end{document}
